\documentclass[11pt,a4paper]{article}

\usepackage{iftex}
\RequirePDFTeX
\usepackage{cmap}
\usepackage[T2A]{fontenc}
\usepackage[utf8]{inputenc}
\usepackage[english,russian]{babel}
\usepackage{PTSerif}
\usepackage{amsmath,amssymb,microtype}
\usepackage[a4paper,left=23mm,right=23mm,top=20mm,bottom=21mm]{geometry}
\usepackage{xcolor,enumitem,fancyhdr,hyperref}

\definecolor{accent}{HTML}{187AA5}
\definecolor{ink}{HTML}{302B28}
\definecolor{muted}{HTML}{6D6560}
\definecolor{tint}{HTML}{EDF4F6}
\definecolor{rulecolor}{HTML}{A9CADA}
\hypersetup{unicode=true,colorlinks=true,urlcolor=accent,linkcolor=accent,
  pdftitle={Домашнее задание 7. Задача выведения КА на орбиту},
  pdfauthor={Чиняев Владимир Николаевич}}
\urlstyle{same}
\color{ink}
\setlength{\parindent}{0pt}
\setlength{\parskip}{3pt}
\linespread{1.08}
\setlist[enumerate]{leftmargin=*,label=\textcolor{accent}{\arabic*.},
  itemsep=3pt,topsep=4pt,parsep=0pt}
\setlist[itemize]{leftmargin=*,itemsep=3pt,topsep=3pt,parsep=0pt}
\setlength{\emergencystretch}{1em}
\pagestyle{fancy}
\fancyhf{}
\renewcommand{\headrulewidth}{0pt}
\renewcommand{\footrulewidth}{0.4pt}
\renewcommand{\footrule}{\hbox to\headwidth{\color{rulecolor}\leaders\hrule height\footrulewidth\hfill}}
\fancyfoot[L]{\footnotesize\color{muted}Введение в маневрирование КА\quad /\quad ДЗ 7}
\fancyfoot[R]{\footnotesize\color{muted}\thepage}

\newcommand{\taskheading}[2]{%
  \par\vspace{5pt}{\large\bfseries\color{accent}Задание #1. #2}\par\smallskip}
\newcommand{\hint}[1]{%
  \par\nobreak\smallskip\noindent
  \begingroup
  \setlength{\fboxsep}{7pt}%
  \setlength{\fboxrule}{0.4pt}%
  \fcolorbox{rulecolor}{tint}{%
    \begin{minipage}{\dimexpr\linewidth-2\fboxsep-2\fboxrule\relax}
    \small\textbf{\color{accent}Подсказка.} #1\par
    \end{minipage}}%
  \endgroup\par\smallskip}

\begin{document}

{\small\color{accent}ВВЕДЕНИЕ В МАНЕВРИРОВАНИЕ КА}\par
\vspace{3pt}
{\huge Домашнее задание 7}\par
{\large Задача выведения КА на орбиту}\par
{\small\color{muted}К лекции 10}
\par{\small\bfseries\color{accent}Срок сдачи: 30 ноября 2026 года, 23:59}
\vspace{4pt}
{\color{rulecolor}\hrule}
\vspace{5pt}

Рассчитайте характеристическую скорость двухступенчатой ракеты
и исследуйте влияние широты космодрома на геометрию запуска.
Оформите решение обоих заданий в предоставленном Jupyter Notebook.

\taskheading{1}{Характеристическая скорость двухступенчатой ракеты}

\begingroup
\setlength{\fboxsep}{9pt}
\colorbox{tint}{\begin{minipage}{\dimexpr\linewidth-2\fboxsep\relax}
\textbf{Исходные данные}
\par\smallskip
Масса полезной нагрузки --- $5$~т. Характеристики ступеней:
\begin{center}
\renewcommand{\arraystretch}{1.12}
\begin{tabular}{lrr}
\hline
Параметр & Первая ступень & Верхняя ступень\\\hline
Масса топлива, т & $160$ & $45$\\
Сухая масса, т & $16$ & $4$\\
Тяга, МН & $3{,}5$ & $0{,}8$\\
Удельный импульс, с & $300$ & $450$\\\hline
\end{tabular}
\end{center}
\vspace{-3pt}
Примите $g_0=9{,}80665$~м/с$^2$. Тяга и удельный импульс постоянны.
Первая ступень расходует всё топливо и отделяется; сразу после этого
включается верхняя ступень, которая также расходует всё топливо.
Начало работы первой ступени соответствует $t=0$.
\end{minipage}}
\endgroup

\begin{enumerate}
\item Найдите массу ракеты при старте, перед разделением,
после разделения и после окончания работы верхней ступени.
\item Рассчитайте длительности работы двигателей,
характеристические скорости обеих ступеней и их сумму.
\item Постройте на общей временной шкале два графика:
массу $m(t)$ и накопленную характеристическую скорость
$\Delta V_{\mathrm{th}}(t)$. Отметьте момент разделения.
Объясните, какая из величин испытывает скачок и почему.
\item Повторите расчёт суммарной характеристической скорости для случая,
когда отработавшая первая ступень \textbf{не отделяется}.
Остальные данные сохраните. Сравните результаты и объясните,
почему отделение повышает возможности ракеты,
хотя само по себе не создаёт скачка скорости.
\end{enumerate}

\hint{Полная масса включает топливо, конструкции обеих ступеней и нагрузку.
Для активного участка используйте $c=g_0I_{\mathrm{sp}}$,
$\dot m=-F/c$ и $\Delta V_{\mathrm{th}}=c\ln(m_b/m_e)$.
Верхняя ступень начинает работу с массы после разделения;
накопленная характеристическая скорость содержит вклад предыдущего участка.
Если первая ступень остаётся, её сухая масса входит и в начальную,
и в конечную массу верхнего активного участка.}

\newpage
{\small\color{muted}Домашнее задание 7\hfill К лекции 10}\par
\vspace{5pt}

\taskheading{2}{Широта космодрома и азимут запуска}

\begingroup
\setlength{\fboxsep}{9pt}
\colorbox{tint}{\begin{minipage}{\dimexpr\linewidth-2\fboxsep\relax}
\textbf{Исходные данные}
\par\smallskip
Требуется прямое выведение на орбиту с наклонением $i=51{,}6^\circ$,
без манёвра изменения плоскости. Рассмотрите условные космодромы:
\[
 \phi=0^\circ,\quad28{,}5^\circ,\quad45^\circ,\quad60^\circ.
\]
Используйте сферическую модель Земли:
\[
 R_E=6378{,}137\ \text{км},\qquad
 \omega_E=7{,}2921159\cdot10^{-5}\ \text{рад/с}.
\]
Азимут $A$ --- азимут \textbf{инерциальной} горизонтальной скорости,
отсчитываемый от севера к востоку.
Оценку выполняйте по геометрии запуска, рассмотренной на лекции.
\end{minipage}}
\endgroup

\begin{enumerate}
\item Для каждого космодрома найдите начальную инерциальную скорость
площадки, обусловленную вращением Земли. Представьте её модуль в м/с
и укажите направление.
\item Из соотношения $\cos i=\cos\phi\sin A$ получите условие существования
геометрически допустимого запуска. Для каждого космодрома найдите
все азимуты $0^\circ\leq A\leq180^\circ$ либо укажите отсутствие решения.
Составьте таблицу результатов.
\item Постройте обе ветви зависимости $A(\phi)$ при
$0^\circ\leq\phi\leq70^\circ$.
Недопустимые широты пропускайте;
отметьте границу области существования решения.
\item Объясните различие двух найденных азимутов.
Какое наклонение получается при запуске строго на восток
и почему оно является минимальным для прямого запуска с данной широты?
\end{enumerate}

\hint{Скорость площадки равна
$V_{E,0}=\lVert\boldsymbol V_0\rVert=\omega_E R_E\cos\phi$
и направлена на восток.
При решении уравнения для азимута сначала проверьте допустимость
значения $\sin A$, затем учтите обе ветви синуса на $[0,\pi]$.
Недопустимые точки графика обозначайте \texttt{np.nan},
не заменяйте их граничным значением.
В граничной допустимой точке две ветви совпадают.}

\vspace{4pt}
\textbf{Работа с ноутбуком и самопроверка}

В \texttt{hw\_7.ipynb} готовы константы, функции для таблиц и графиков,
чекеры \texttt{check\_task1} и \texttt{check\_task2}.
Заполните отмеченные переменные, расчётные функции и текстовые ячейки.
Использование чекеров и сохранение результатов их вызова не требуются.
В коде используйте килограммы, ньютоны, секунды, метры и радианы;
в таблицах и на графиках массы приводите в тоннах, углы --- в градусах.
Не округляйте промежуточные результаты.

\vspace{5pt}
{\color{rulecolor}\hrule}
\vspace{5pt}
\textbf{Что представить}

Сдайте один ноутбук с решениями обоих заданий: расчётными формулами,
таблицами, графиками и объяснениями.
Перед отправкой перезапустите ядро, выполните ячейки по порядку
и сохраните ноутбук с результатами.
Название файла: \texttt{hw\_7\_name\_surname.ipynb};
вместо \texttt{name\_surname} укажите имя и фамилию латиницей.

\end{document}
